\clearpage
\phantomsection
\addcontentsline{toc}{section}{Introducción}

{\centering\Large\bfseries INTRODUCCIÓN\par}
\vspace{0.8cm}

En el entorno empresarial contemporáneo, la administración de los empleados y el procesamiento de compensaciones constituyen procesos críticos para la continuidad operativa. No obstante, una gran parte de las organizaciones continúa gestionando sus planillas mediante hojas de cálculo convencionales. Este enfoque tradicional demanda una inversión desmedida de tiempo en la validación manual de asistencias, permisos, bonificaciones y deducciones de ley, lo que incrementa sustancialmente el margen de error humano en el cálculo de la planilla, generando retrasos en la dispersión de fondos, pagos inexactos y potenciales contingencias laborales. Asimismo, el almacenamiento de información financiera en archivos planos o locales expone a las organizaciones a graves riesgos de fuga de datos o accesos no autorizados.

Frente a este escenario, el presente proyecto de innovación propone el diseño e implementación de un sistema web integral bajo una arquitectura de despliegue monolítica modular, orientado a centralizar la gestión de personal y automatizar el cálculo sistemático de planillas. Para optimizar el ciclo operativo, la plataforma se integra con la herramienta de orquestación y automatización de flujos \textbf{n8n}, la cual procesa la dispersión de pagos y gestiona el envío automatizado de las boletas de pago digitales a las casillas de correo electrónico de cada colaborador.

Con el fin de garantizar la confidencialidad, integridad y disponibilidad de los datos sensibles, el sistema incorpora un modelo riguroso de control de acceso basado en roles (RBAC), mecanismos de cifrado \textbf{AES-256-GCM} para los registros financieros sensibles y un módulo de auditoría continua que registra la trazabilidad completa de cada transacción y cambio de estado ejecutado en la plataforma.